% titel: Blatt 9, A7: MWG und Gleichgewichtsverschiebung
% Teil II: der durchgerechnete Weg.

\subsection*{Reaktion}
$2H_2S(g) + CH_4(g) \rightleftharpoons CS_2(g) + 4H_2(g)$, von links nach
rechts \textbf{exotherm}.

\subsection*{a) K$_c$}
\[
  K_c = \frac{[CS_2][H_2]^4}{[H_2S]^2[CH_4]}.
\]

\subsection*{b) K$_p$ und Zusammenhang}
\[
  K_p = \frac{p_{CS_2}\,p_{H_2}^4}{p_{H_2S}^2\,p_{CH_4}},
  \qquad
  K_p = K_c \,(RT)^{\Delta n},\quad
  \Delta n = (1+4)-(2+1) = +2 .
\]

\subsection*{c) Gleichgewichtsverschiebung}
\begin{itemize}
  \item \textbf{Temperaturerhöhung.} Die Vorwärtsreaktion ist exotherm; nach
        Le Chatelier begünstigt Wärmezufuhr die endotherme Rückreaktion
        $\Rightarrow$ Gleichgewicht verschiebt nach \emph{links}.
  \item \textbf{Druckerhöhung.} Links 3, rechts 5 Gaspartikel; Druckerhöhung
        begünstigt die \emph{kleinere} Teilchenzahl $\Rightarrow$ nach
        \emph{links}.
  \item \textbf{Zugabe von H$_2$S.} Erhöht die Edukt-Konzentration; das
        Gleichgewicht verschiebt, um zu kompensieren, nach \emph{rechts}
        (mehr CS$_2$ und H$_2$).
  \item \textbf{Katalysator.} Beschleunigt Hin- und Rückreaktion gleicher
        massen; die \emph{Lage} des Gleichgewichts ändert sich \emph{nicht}.
\end{itemize}

\begin{ergebnis}
$K_c$ und $K_p$ wie oben; $K_p=K_c(RT)^2$. Temperaturerhöhung und
Druckerhöhung $\Rightarrow$ nach links; H$_2$S-Zugabe $\Rightarrow$ nach
rechts; Katalysator $\Rightarrow$ keine Lageänderung.
\end{ergebnis}

\begin{kontrolle}
Konsistenz: exotherme Vorwärtsrichtung und mehr Teilchen rechts machen
beide Druck- und Temperaturantworten plausibel. Der Katalysator-Teil wird
durch die Unabhängigkeit von $K_c$ von der Beschleunigung bestätigt.
\end{kontrolle}

\begin{merksatz}
Kp aus Partialdrücken, $K_p=K_c(RT)^{\Delta n}$. Druck begünstigt weniger
Teilchen, Wärme begünstigt die endotherme Richtung, Edukt-Zugabe verschiebt
zur Produktseite, Katalysator ändert die Lage nie.
\end{merksatz}
